%----------------------------------------------------------------------
\section{SPPF: the paper as a discrimination regime}
\label{app:sppf}
%----------------------------------------------------------------------

This appendix presents the paper's dependency structure as
a Shared Packed Parse Forest (SPPF), treating every formal
object --- definition, theorem, proposition, corollary,
conjecture, example, and remark --- as a node in a
discrimination regime. The SPPF is the framework pointed
at itself.

\subsection{The meta-discrimination}

The paper is a population of 125 formal objects. Four
discriminators interrogate this population:

\begin{center}
\small
\begin{tabular}{c|l|l}
\textbf{Tag} & \textbf{Discriminator} & \textbf{Fires when} \\
\hline
\textsf{A} & Algebraic & Always (every node uses
  $\GF(2)$, $\bigwedge V$) \\
\textsf{P} & Profile linearity &
  Node depends on Def \ref{def:profile-linear} \\
\textsf{E} & Evaluation linearity &
  Node depends on Def \ref{def:eval-linear} \\
\textsf{O} & Operationalist axiom &
  Node depends on Def \ref{def:operationalist}
\end{tabular}
\end{center}

The \emph{meta-membership-profile} of each node is its
axiom class: the subset $\{\textsf{A}, \textsf{P},
\textsf{E}, \textsf{O}\}$ of discriminators that fire.
A node with class \textsf{A} depends only on the algebra.
A node with class \textsf{APE} depends on both regulatory
axioms. The profile is a composed discriminator
(Definition \ref{def:profile}): the pair of meta-$\tau$
(less-refined axiom set) and meta-$\sigma$ (more-refined
axiom set) is the Grothendieck construction applied to
the paper's own morphisms.

The meta-$\kappa$-regime evolves as the paper proceeds:
\S2 has zero nodes, \S3 introduces the algebraic
vocabulary (all \textsf{A}), \S4 introduces the
regulatory axioms (\textsf{P}, \textsf{E}) and the
epistemic axiom (\textsf{O}), and subsequent sections
inherit their dependencies transitively. The SPPF is the
$\kappa$-evolution trace of the paper's own construction.

\subsection{Axiom class summary}

\begin{center}
\begin{tabular}{l|r|l}
\textbf{Class} & \textbf{Count} & \textbf{Interpretation} \\
\hline
\textsf{A} & 101 & Algebra alone (81\%) \\
\textsf{AP} & 9 & Algebra + profile linearity \\
\textsf{AEP} & 6 & Algebra + both linearities \\
\textsf{AO} & 8 & Algebra + operationalist axiom \\
\textsf{AOP} & 1 & All three axiom types \\
\hline
Total & 125 &
\end{tabular}
\end{center}

81\% of the paper's formal objects depend only on the
algebra (as stated in the reading guide,
\S\ref{sec:reading-guide}). Profile linearity adds 9, evaluation linearity
adds 6 (always paired with profile linearity), and the
operationalist axiom adds 8. Only one node
(Remark 4.10, general pathology protection) depends on
all three regulatory/epistemic commitments simultaneously.

\subsection{Complete node inventory}

Each row gives the node's section, line, kind
(\textsf{D}ef, \textsf{T}hm, \textsf{P}rop,
\textsf{C}or, \textsf{J}conj, \textsf{E}x,
\textsf{R}em), title, and axiom class.


\medskip
\noindent\textbf{\S01 Introduction} (1 nodes).

\begin{center}
\scriptsize
\begin{tabular}{r|c|l|c}
\textbf{Line} & \textbf{Kind} & \textbf{Title} & \textbf{Class} \\
\hline
36 & \textsf{D} & Operationalist axiom & \textbf{AO} \\
\end{tabular}
\end{center}

\medskip
\noindent\textbf{\S03 Framework} (17 nodes).

\begin{center}
\scriptsize
\begin{tabular}{r|c|l|c}
\textbf{Line} & \textbf{Kind} & \textbf{Title} & \textbf{Class} \\
\hline
7 & \textsf{D} & Discriminator & A \\
17 & \textsf{D} & $\tau$ and $\sigma$ as datum properties & A \\
32 & \textsf{D} & Membership profile & A \\
60 & \textsf{R} & The four-cell structure & A \\
114 & \textsf{D} & Profile linearity & \textbf{AP} \\
146 & \textsf{D} & XOR signature & A \\
158 & \textsf{D} & Residue & A \\
169 & \textsf{R} & Two origins of residue & A \\
182 & \textsf{D} & Discrimination regime & A \\
192 & \textsf{D} & $\kappa$-evolution & A \\
199 & \textsf{P} & Monotonicity of $\kappa$ & A \\
212 & \textsf{D} & Write path & A \\
218 & \textsf{D} & Read path & A \\
227 & \textsf{P} & Order-independence of filtration & A \\
244 & \textsf{C} & Well-definedness of type erasure & A \\
251 & \textsf{D} & Membership & A \\
261 & \textsf{R} & Yoneda content & A \\
\end{tabular}
\end{center}

\medskip
\noindent\textbf{\S04 Set theory} (25 nodes).

\begin{center}
\scriptsize
\begin{tabular}{r|c|l|c}
\textbf{Line} & \textbf{Kind} & \textbf{Title} & \textbf{Class} \\
\hline
13 & \textsf{D} & $\kappa$-equivalence & A \\
21 & \textsf{T} & Extensionality & A \\
46 & \textsf{D} & Discriminative comprehension & \textbf{AP} \\
59 & \textsf{T} & Exclusion of Russell's predicate & \textbf{AEP} \\
100 & \textsf{D} & Evaluation linearity & \textbf{AEP} \\
179 & \textsf{R} & Profile linearity as the load-bearing axiom & \textbf{AEP} \\
191 & \textsf{R} & The restriction is relocated & \textbf{AEP} \\
217 & \textsf{R} & General pathology protection & \textbf{AOP} \\
246 & \textsf{D} & Boolean dimension & A \\
257 & \textsf{P} & Non-default Booleanness & A \\
283 & \textsf{T} & Partial foundation & \textbf{AEP} \\
308 & \textsf{R} & Cycles of length $\geq 2$ are permitted & \textbf{AP} \\
374 & \textsf{P} & Empty set & A \\
383 & \textsf{P} & Pairing & \textbf{AP} \\
395 & \textsf{P} & Symmetric difference as primitive & A \\
421 & \textsf{D} & Discriminative power set & A \\
436 & \textsf{J} & Constructive continuum & A \\
462 & \textsf{D} & Discriminative infinity & A \\
530 & \textsf{R} & Foundational status of indistinguishability & \textbf{AO} \\
555 & \textsf{R} & Epistemic asymmetry and the halting problem & \textbf{AO} \\
578 & \textsf{R} & The hyper-recursive tower & \textbf{AO} \\
621 & \textsf{R} & The proof assistant concession & \textbf{AO} \\
681 & \textsf{T} & Choice from residue consumption & \textbf{AO} \\
807 & \textsf{R} & Choice as a dependent type via Diaconescu & A \\
840 & \textsf{E} & Banach--Tarski is non-paradoxical & A \\
\end{tabular}
\end{center}

\medskip
\noindent\textbf{\S05 Homotopy} (11 nodes).

\begin{center}
\scriptsize
\begin{tabular}{r|c|l|c}
\textbf{Line} & \textbf{Kind} & \textbf{Title} & \textbf{Class} \\
\hline
15 & \textsf{D} & Degeneracy as identity & A \\
40 & \textsf{D} & Discrimination complex & A \\
64 & \textsf{D} & Boundary operator & A \\
82 & \textsf{P} & Chain complex & A \\
138 & \textsf{T} & Exhaustivity is operational & A \\
184 & \textsf{C} & Non-uniqueness under truncation & A \\
240 & \textsf{T} & Intrinsic groupoid structure & A \\
308 & \textsf{P} & Univalence as extensionality & A \\
338 & \textsf{R} & Transport as reclassification & A \\
357 & \textsf{D} & Emergent $n$-groupoid & A \\
375 & \textsf{P} & Growth of homotopical complexity & A \\
\end{tabular}
\end{center}

\medskip
\noindent\textbf{\S06 Categories} (19 nodes).

\begin{center}
\scriptsize
\begin{tabular}{r|c|l|c}
\textbf{Line} & \textbf{Kind} & \textbf{Title} & \textbf{Class} \\
\hline
16 & \textsf{R} & The source of asymmetry & A \\
27 & \textsf{D} & Directed morphism & A \\
40 & \textsf{P} & Non-invertibility & A \\
60 & \textsf{C} & Cost of isomorphism & A \\
70 & \textsf{D} & Emergent category & A \\
91 & \textsf{D} & Composition via filtration & A \\
107 & \textsf{R} & The framework is 2-categorical from the start & A \\
129 & \textsf{T} & 2-category axioms & A \\
162 & \textsf{D} & Functor & A \\
178 & \textsf{P} & Forgetful and free functors & A \\
241 & \textsf{T} & Triadic adjunction system & A \\
346 & \textsf{C} & The free-forgetful adjunction is one of six & A \\
381 & \textsf{D} & Natural transformation & A \\
399 & \textsf{P} & Naturality from coherence of discrimination & A \\
439 & \textsf{D} & Representable presheaf & A \\
447 & \textsf{T} & Yoneda & A \\
510 & \textsf{D} & Limit as universal cone & A \\
523 & \textsf{P} & Limits from exhaustive filtration & A \\
560 & \textsf{R} & Completeness is not a fixed property & A \\
\end{tabular}
\end{center}

\medskip
\noindent\textbf{\S07 Higher structure} (16 nodes).

\begin{center}
\scriptsize
\begin{tabular}{r|c|l|c}
\textbf{Line} & \textbf{Kind} & \textbf{Title} & \textbf{Class} \\
\hline
14 & \textsf{T} & Natural $(n,1)$-structure & A \\
70 & \textsf{D} & Directed homotopy & A \\
117 & \textsf{E} & Directed 2-morphism: two classification schemes & A \\
158 & \textsf{P} & Structure of $k$-morphisms & A \\
173 & \textsf{R} & Truncation level as complexity measure & A \\
196 & \textsf{D} & Triangle identity & A \\
216 & \textsf{R} & Grounding the identification & A \\
254 & \textsf{P} & Three perspectives & A \\
289 & \textsf{P} & Dynamics transfer under permutation & A \\
334 & \textsf{R} & Semantic asymmetry across perspectives & \textbf{AEP} \\
350 & \textsf{D} & $\GF(2)$ toroid & A \\
376 & \textsf{T} & Free $(n,k)$-structure & A \\
424 & \textsf{C} & Direction requires no additional axiom & A \\
433 & \textsf{R} & Geometric content of the toroid & A \\
451 & \textsf{R} & Why $(\infty,1)$-categories appear natural & A \\
475 & \textsf{R} & The $(\infty,\infty)$-category & \textbf{AO} \\
\end{tabular}
\end{center}

\medskip
\noindent\textbf{\S08 Monoidal and closure} (13 nodes).

\begin{center}
\scriptsize
\begin{tabular}{r|c|l|c}
\textbf{Line} & \textbf{Kind} & \textbf{Title} & \textbf{Class} \\
\hline
4 & \textsf{D} & Monoidal product & A \\
19 & \textsf{P} & Symmetric monoidal structure & A \\
59 & \textsf{T} & Delooping equivalence & A \\
112 & \textsf{C} & Delooping collapse over $\GF(2)$ & A \\
159 & \textsf{D} & Swap conjugation & A \\
169 & \textsf{D} & Cayley--Dickson tower over $\GF(2)^2$ & A \\
199 & \textsf{J} & The periodic table is two-dimensional & A \\
249 & \textsf{R} & The swap conjugation is already present & A \\
281 & \textsf{R} & Open verification: hexagon axioms at CD step 2 & A \\
306 & \textsf{D} & Internal hom & A \\
316 & \textsf{T} & Monoidal closure & A \\
362 & \textsf{R} & Closure is the triangle identity & A \\
374 & \textsf{D} & Evaluation morphism & A \\
\end{tabular}
\end{center}

\medskip
\noindent\textbf{\S09 Topos and self-reference} (22 nodes).

\begin{center}
\scriptsize
\begin{tabular}{r|c|l|c}
\textbf{Line} & \textbf{Kind} & \textbf{Title} & \textbf{Class} \\
\hline
15 & \textsf{D} & Enrichment data & A \\
23 & \textsf{T} & Canonical self-enrichment & A \\
43 & \textsf{D} & Subobject classifier & A \\
48 & \textsf{T} & $\Omega$ is a subobject classifier & A \\
76 & \textsf{D} & The fibered category & A \\
84 & \textsf{T} & Fibered topos & \textbf{AP} \\
102 & \textsf{C} & Pro-system interpretation & \textbf{AO} \\
112 & \textsf{D} & Evidence semantics & A \\
118 & \textsf{D} & Information ordering & A \\
125 & \textsf{P} & Residues are extremal & A \\
132 & \textsf{T} & Dynamic proof theory & A \\
151 & \textsf{P} & First-order logic is internal & \textbf{AP} \\
166 & \textsf{P} & Boolean logic as a dependent type & A \\
177 & \textsf{T} & Self-hosting & \textbf{AP} \\
191 & \textsf{T} & Discrimination is a fixed point & A \\
203 & \textsf{T} & Discrimination is irremovable & \textbf{AP} \\
216 & \textsf{D} & The site of regimes & A \\
223 & \textsf{T} & Sheaf topos & \textbf{AP} \\
242 & \textsf{D} & Coverage profile & A \\
249 & \textsf{T} & Fano closure & A \\
259 & \textsf{C} & Self-modeling & A \\
266 & \textsf{T} & Convergence rate & A \\
\end{tabular}
\end{center}

\medskip
\noindent\textbf{\S10 Discussion} (1 nodes).

\begin{center}
\scriptsize
\begin{tabular}{r|c|l|c}
\textbf{Line} & \textbf{Kind} & \textbf{Title} & \textbf{Class} \\
\hline
270 & \textsf{R} & Epistemic limitation of self-audit & A \\
\end{tabular}
\end{center}

\subsection{Critical dependency chains}

The longest chains, with axiom class of each link:

\emph{Russell exclusion} (length 3, class \textsf{AEP}):
Def 3.1 [\textsf{A}] $\to$ Def 3.5 [\textsf{AP}]
$\to$ Def 4.6 [\textsf{AEP}] $\to$ Thm 4.5 [\textsf{AEP}].

\emph{Fibered topos} (length 7, class \textsf{AP}):
Def 3.1 [\textsf{A}] $\to$ Def 5.2 [\textsf{A}]
$\to$ Prop 5.4 [\textsf{A}] $\to$ Thm 5.7 [\textsf{A}]
$\to$ Thm 8.13 [\textsf{A}] $\to$ Def 9.4 [\textsf{A}]
$\to$ Thm 9.6 [\textsf{AP}].

\emph{Infinite choice} (length 7, class \textsf{AO}):
Def 3.1 [\textsf{A}] $\to$ Def 3.7 [\textsf{A}]
$\to$ Def 3.9 [\textsf{A}] $\to$ Thm 5.7 [\textsf{A}]
$\to$ Prop 4.26 [\textsf{A}] $\to$ Def 1.1 [\textsf{AO}]
$\to$ Thm 4.31 [\textsf{AO}].

\emph{Irremovability} (length 8+, class \textsf{AP}):
$\cdots \to$ Thm 5.7 [\textsf{A}] $\to$ Thm 6.11
[\textsf{A}] $\to$ Thm 9.6 [\textsf{AP}] $\to$
Thm 9.14 [\textsf{AP}]. Uses the framework's own
machinery; circularity acknowledged (Remark 10.1).

\subsection{The SPPF as successive refinement}

The SPPF is a magma under predicate accumulation: the
product of two nodes is their composition (the wedge
product of their dependency sets). The axiom class of a
composed result is the union of the constituent classes.

Reading the paper sequentially, each section refines the
previous sections' results by adding new predicates
(discriminators at the meta-level):

\begin{itemize}[nosep]
  \item \S3: introduces the algebraic vocabulary. All
    nodes are class \textsf{A}. Meta-$\kappa$ has one
    dimension.
  \item \S4: introduces \textsf{P}, \textsf{E},
    \textsf{O}. Meta-$\kappa$ evolves to four dimensions.
    The Russell exclusion is the first node requiring all
    three regulatory/epistemic dimensions.
  \item \S\S5--8: no new axioms. These sections refine
    the \textsf{A}-class nodes by accumulating algebraic
    structure. The meta-$\kappa$ does not evolve (no new
    axioms), but the population grows.
  \item \S9: inherits \textsf{P} from the topos
    construction. The meta-profile of \S9's theorems
    records which axioms the topos transitively requires.
\end{itemize}

At the meta-level, $\tau$ is the less-refined set of
predicates (the axiom class inherited from dependencies)
and $\sigma$ is the more-refined set (the specific
combination of axioms the node directly invokes). A node
that inherits \textsf{AP} from its dependencies but does
not directly invoke profile linearity has meta-$\tau$
active and meta-$\sigma$ silent on the \textsf{P}
dimension: the dependency is transitive, not direct. This
distinction is the meta-level residue that the SPPF tracks.

The SPPF is the Grothendieck construction applied to the
paper's own morphisms: it fibers the population of formal
objects over the poset of axiom classes, with the
transition functors being the dependency edges. The
resulting structure is the paper viewed as its own topos.
